
\documentclass[12pt,a4paper]{article}

% Поддержка русского языка
\usepackage[T2A]{fontenc}
\usepackage[utf8]{inputenc}
\usepackage[russian]{babel}

% Математика
\usepackage{amsmath,amssymb,amsthm}

% Настройка страницы
\usepackage[a4paper,margin=2.5cm]{geometry}

\begin{document}
	
	\tableofcontents
	
	\newpage
	\section{Lagrange polynomials}
	
	\subsection{Interpolation}
	Пусть $(x_k, y_k)_{k=0}^{n}$ точки из $\mathbb{R}^2$ \\
	Задача: найти $p \in \mathbb{P}_n$ т.ч. $p(x_k)=y_k, k=0 \ldots n$,\\
	где $\mathbb{P}_n$ -- многочлены степени не выше $n$\\
	
	\noindent Введём "базисные" многочлены:\\
	\[
	L_i(x) := \underset{i \neq j}{\prod_{j=0}^{n}} \frac{x - x_j}{x_i - x_j}, \ i \in 0 \ldots n
	\]
	
	\noindent Заметим, что:
	\[
	L_i(x_k) = 
	\begin{cases}
	1, \ i=k \\ 
	0, \ i \neq k\\
	\end{cases}i,k \in 0 \ldots n
	\]
	
	\noindent Введём интерполяционный многочлен:
	\[
	L(x) := \sum_{k=0}^{n} y_k L_k(x)
	\]
	
	
	\noindent Очевидно, это то что надо
	\[
	L(x_k) = y_k, \ k \in 0 \ldots n
	\]
	
	\noindent Таким образом, мы нашли базис $\mathbb{P}_n$ т.к. $\forall q \in \mathbb{P}_n, \ \forall(x_0 \ldots x_n)$
	\[
	q(x_k) = \sum_{k=0}^{n} q(x_k) L_k(x_k), \ k \in 0 \ldots n
	\]
	
	\noindent но разность левой и правой части это многочлен степени не выше $n$, а мы уже нашли $n+1$ корень. Значит они тождественно равны:
	\[
	q(x) = \sum_{k=0}^{n} q(x_k) L_k(x), \ \forall x \in \mathbb{R}
	\]
	
	\newpage
	\subsection{Remainder}
	Пусть $p \in \mathbb{P}_n$\\
	Пусть $(x_k, y_k)_{k=1}^{n}$ точки из $\mathbb{R}^2$ \\
	Можно построить
	\[
	L(x) = \sum_{k=1}^{n} p(x_k) L_k(x)
	\]
	
	\noindent Рассмотрим остаток
	\[
	R(x) = p(x) - L(x)
	\]
	
	\noindent Это многочлен степени не выше $n$ и мы знаем что $x_k, k \in 1 \ldots n$ это его корни. Также из того что $L_k$ строились по n узлам мы имеем $deg(L) \leq n-1$, 
	то есть старший коэффициент $a_n(R)$ многочлена $R$ совпадает со старшим коэфом $a_n(p)$ многочлена $p$.
	Итоговый вид остатка:
	\[
	R(x) = a_n(p) \prod_{k=1}^{n} (x-x_k) \\
	\]
	
	\noindent Представление для $p$:
	\begin{equation}
	p(x) = 
	L(x) + R(x) = 
    \sum_{k=1}^{n} p(x_k) L_k(x) + 
    a_n(p) \prod_{k=1}^{n} (x-x_k)
    \label{eq:lagr_expansion}
	\end{equation}
	
	\newpage
	
	\section{Обозначения}
	Фиксируем $x \in \mathbb{R}$\\
	Пусть $f(x) = (1, x, \ldots, x^m)^\mathrm{T}$, $c=f'(x)$\\
	
	\noindent Рассмотрим план
	\[
		\xi = 
		\left(
		\begin{matrix}
		x_1 \ x_2 \ldots x_n \\
		\omega_1 \ \omega_2 \ldots \omega_n	
		\end{matrix}
		\right)
	\]
	
	\noindent Та самая система для плана с $n$ узлами имеет вид:
	\[
	c = h M(\xi) p = h \sum_{k=1}^{n} \omega_k <p,f(x_k)> f(x_k)
	\]	
	
	\noindent или, c учётом $h = (\lambda ^*)^{-1}$ и чередования знаков т.е. $<p,f(x_k)> = (-1)^{k+1}$
	\[
	\lambda ^*
	\begin{pmatrix}
	0 \\
	1 \\
	\vdots \\
	m x^{m-1}
	\end{pmatrix} = 
	\sum_{k=1}^{n} a_k 
	\begin{pmatrix}
		1 \\
		x_k \\
		\vdots \\
		x_k^{m}
	\end{pmatrix}, \quad a_k = \omega_k <p,f(x_k)>
	\]
	
	\noindent Введём линейные операторы $A_1, A_2, A: \mathbb{P}_m \to \mathbb{R}$:
	\[
	A_1(p) := \lambda^* p'(x)
	\]
	\[
	A_2(p) := \sum_{k=1}^{n} a_k p(x_k)
	\]
	\[
	A(p) := A_1(p) - A_2(p)
	\]
	
	\noindent Нетрудно видеть, что наша система эквивалентна вот такому:
	\[
	\lbrace 1, x, \ldots, x^m \rbrace \subset \ker A
	\]
	
	\noindent Т.к $A$ линейно, то это в свою очередь эквивалентно
	\begin{equation}
	\mathbb{P}_m \subset \ker A
	\label{eq:cond1}
	\end{equation}
	
	
	\newpage
	\section{Чебышёвские планы}
	\subsection{Ищем веса в чебышёвских планах}
	
	
	\noindent (\textbf{в этом разделе} $\mathbf{n = m+1}$)\\ 
	
	
	\noindent Вспоминаем, что \eqref{eq:lagr_expansion}  $\forall p \in \mathbb{P}_m$ даёт нам
	\[
	p(t) = 
	\sum_{k=1}^{m+1} p(x_k) L_k(t), \  \forall t \in \mathbb{R}
	\]
	
	\noindent Дифференцируя это равенство по $t$ и подставляя $t=x$ получаем 
	\[
	p'(x) = 
	\sum_{k=1}^{m+1} p(x_k) L_k'(x)
	\]
	
	
	\noindent Получили, что $\forall p \in \mathbb{P}_m$
	\begin{equation}
		\lambda^* p'(x) = 
		\sum_{k=1}^{m+1} p(x_k) ( \lambda^* L_k'(x) ) 
		\label{eq:re1_2}
	\end{equation}
	
	\noindent Но из условия \eqref{eq:cond1} 
	\[
	A(p) = 0, \quad \forall p \in \mathbb{P}_m
	\]
	
	\noindent То есть $\forall p \in \mathbb{P}_m$
	\begin{equation}
		\lambda^* p'(x) = 
		\sum_{k=1}^{m+1} p(x_k) ( a_k ) 
		\label{eq:re2_2}
	\end{equation}
	
	\noindent Из \eqref{eq:re1_2} и \eqref{eq:re2_2} получаем $\forall p \in \mathbb{P}_m$
	\[
	\sum_{k=1}^{m+1} p(x_k) ( a_k ) = \sum_{k=1}^{m+1} p(x_k) ( \lambda^* L_k'(x) ) 
	\] 
	
	\noindent Выбирая в качестве $p$ штуки которые равны $1$ и $0$ где надо ($L_k(x) ???$), делаем вывод
	\[
	a_k = \lambda^* L_k'(x), \ \forall k \in 1 \ldots m+1 
	\]
	\[
	\omega_k (-1)^{k+1} = \lambda^* L_k'(x), \ \forall k \in 1 \ldots m+1 
	\]
	\[
	\omega_k  = (-1)^{k+1} \lambda^* L_k'(x), \ \forall k \in 1 \ldots m+1 
	\]
	
	
	\noindent Юзая условие на веса $\sum_{k=1}^{m+1} \omega_k = 1$ ищем лямбду:
	\[
	1 = \sum_{k=1}^{m+1} \omega_k = \sum_{k=1}^{m+1} (-1)^{k+1} \lambda^* L_k'(x) = \lambda^* \sum_{k=1}^{m+1} (-1)^{k+1}  L_k'(x)
	\]
	
	\noindent получаем
	\[
	\lambda^* = \left( \sum_{k=1}^{m+1} (-1)^{k+1}  L_k'(x) \right) ^ {-1}
	\]
	
	\noindent финальное выражение для весов
	\[
	\omega_k  = 
	\frac
	{
		(-1)^{k+1} L_k'(x)
	}
	{
		\sum_{k=1}^{m+1} (-1)^{k+1}  L_k'(x)
	}, \ \forall k \in 1 \ldots m+1 
	\]
	
	\[
	L_k(x) = \underset{k \neq j}{\prod_{j=1}^{m+1}} \frac{x - x_j}{x_k - x_j}, \ k \in 1 \ldots m+1
	\]
	
	\subsection{Ограничение на $x$ от весов}
	
	Мы требуем $\forall k \in 1 \ldots m+1 \ \omega_k > 0 $
	
	\noindent То есть нужно решить систему неравенств относительно $x$:
	\[
	\omega_k  = 
	\frac
	{
		(-1)^{k+1} L_k'(x)
	}
	{
		\sum_{k=1}^{m+1} (-1)^{k+1}  L_k'(x)
	} > 0 , \ \forall k \in 1 \ldots m+1 
	\]
	
	
	
	\subsection{Ищем узлы в чебышёвских планах}
	
	\noindent А узлы это просто точки альтернанса полинома Чебышёва 1-го рода 
	\[
	T_{m}(x) = \cos( m \arccos(x))
	\]
	\noindent Они имеют вид 
	\[
	x_k = \cos \left( \frac{(m-k) \pi}{m} \right), \ k \in 0 \ldots m
	\]
	
	
	
	\newpage
	\section{Получебышёвские планы}
	\subsection{Ищем веса в получебышёвских планах}
	\noindent (\textbf{А теперь} $\mathbf{n = m}$)\\ 
	Теперь рассмотрим 
	\[
	p_0(x) = \prod_{k=1}^{m} (x-x_k)
	\]
	
	\noindent Для него
	\[A(p_0) = 0\]
	
	\noindent И, также 
	\[A_2(p_0) = 0\]
	
	\noindent Отсюда 
	\[A_1(p_0) = \lambda^* p_0'(x) =  0\]
	
	\noindent При фиксированном $x$ получаем условие на $x_k, k \in 1 \ldots m$ (лямбда не ноль)
	\[
	p_0'(x) = \frac{d}{dt} \left[\prod_{k=1}^{m} (t - x_k) \right]_{t=x} = \ 0
	\]
	
	\noindent Вспоминаем, что \eqref{eq:lagr_expansion}  $\forall p \in \mathbb{P}_m$ даёт нам
	\[
	p(t) = 
	\sum_{k=1}^{m} p(x_k) L_k(t) + 
	a_m(p) \prod_{k=1}^{m} (t-x_k) = 
	\sum_{k=1}^{m} p(x_k) L_k(t) + 
	a_m(p) p_0(t)
	\]
	
	\noindent Дифференцируя это равенство по $t$ получаем 
	\[
	p'(t) = 
	\sum_{k=1}^{m} p(x_k) L_k'(t) + 
	a_m(p) p_0'(t)
	\]
	
	\noindent Подставляя $t=x$ убиваем остаток
	\[
	p'(x) = 
	\sum_{k=1}^{m} p(x_k) L_k'(x)
	\]
	
	\noindent Получили, что $\forall p \in \mathbb{P}_m$
	\begin{equation}
	\lambda^* p'(x) = 
	\sum_{k=1}^{m} p(x_k) ( \lambda^* L_k'(x) ) 
	\label{eq:re1}
	\end{equation}
	
	\noindent Но из условия \eqref{eq:cond1} 
	\[
	A(p) = 0, \quad \forall p \in \mathbb{P}_m
	\]
	
	\noindent То есть $\forall p \in \mathbb{P}_m$
	\begin{equation}
	\lambda^* p'(x) = 
	\sum_{k=1}^{m} p(x_k) ( a_k ) 
	\label{eq:re2}
	\end{equation}
	
	\noindent Из \eqref{eq:re1} и \eqref{eq:re2} получаем $\forall p \in \mathbb{P}_m$
	\[
	\sum_{k=1}^{m} p(x_k) ( a_k ) = \sum_{k=1}^{m} p(x_k) ( \lambda^* L_k'(x) ) 
	\] 
	
	\noindent Выбирая в качестве $p$ штуки которые равны $1$ и $0$ где надо, делаем вывод
	\[
	a_k = \lambda^* L_k'(x), \ \forall k \in 1 \ldots m 
	\]
	\[
	\omega_k (-1)^{k+1} = \lambda^* L_k'(x), \ \forall k \in 1 \ldots m 
	\]
	\[
	\omega_k  = (-1)^{k+1} \lambda^* L_k'(x), \ \forall k \in 1 \ldots m 
	\]
	
	
	\noindent Юзая условие на веса $\sum_{k=1}^{m} \omega_k = 1$ ищем лямбду:
	\[
	1 = \sum_{k=1}^{m} \omega_k = \sum_{k=1}^{m} (-1)^{k+1} \lambda^* L_k'(x) = \lambda^* \sum_{k=1}^{m} (-1)^{k+1}  L_k'(x)
	\]
	
	\noindent получаем
	\[
	\lambda^* = \left( \sum_{k=1}^{m} (-1)^{k+1}  L_k'(x) \right) ^ {-1}
	\]
	
	\noindent финальное выражение для весов
	\[
	\omega_k  = 
	\frac
	{
		(-1)^{k+1} L_k'(x)
	}
	{
		\sum_{k=1}^{m} (-1)^{k+1}  L_k'(x)
	}, \ \forall k \in 1 \ldots m 
	\]
	
	\[
	L_k(x) = \underset{k \neq j}{\prod_{j=1}^{m}} \frac{x - x_j}{x_k - x_j}, \ k \in 1 \ldots m
	\]
	
	\subsection{Ограничение на $x$ от весов}
	
	Мы требуем $\forall k \in 1 \ldots m \ \omega_k > 0 $
	
	\noindent То есть нужно решить систему неравенств относительно $x$:
	\[
	\omega_k  = 
	\frac
	{
		(-1)^{k+1} L_k'(x)
	}
	{
		\sum_{k=1}^{m} (-1)^{k+1}  L_k'(x)
	} > 0 , \ \forall k \in 1 \ldots m 
	\]
	
	\newpage
	\subsection{А теперь ищем узлы}
	
	\subsection*{Если растяжка вправо}
	
	\noindent В таком сетапе считаем
	\[
	    -1 = x_1 < x_2 < \ldots < x_m < 1 < z
	\]
	
	\noindent где $x_k = a t_{k-1} + b, k \in 1, \ldots m$ \\
	и $z=a t_m + b$ (её игнорим)\\
	для неких $a, b \in \mathbb{R}$, где $a > 0$ (мы их не знаем, но они явно и не нужны, главное что какое-то линейное преобразование) \\
	а $t_k, \  k \in 0, \ldots, m$ есть точки альтернанса полинома Чебышёва 1-го рода 
	\[
	T_{m}(x) = \cos( m \arccos(x))
	\]
	\noindent Они имеют вид 
	\[
	t_k = \cos \left( \frac{ (m-k) \pi}{m} \right), \ k \in 0 \ldots m
	\]
	
	\noindent Уравнения описывающие растяжение (отношение длин интервалов между точками альтернанса сохраняется)
	\begin{align*}
		\frac{x_3 - x_2}{x_2 - x_1} &= \frac{t_3 - t_2}{t_2 - t_1}\\ 
		\frac{x_4 - x_3}{x_2 - x_1} &= \frac{t_4 - t_3}{t_2 - t_1}\\
		&\cdots\\
		\frac{x_m - x_{m-1}}{x_2 - x_1} &= \frac{t_m - t_{m-1}}{t_2 - t_1}\\
	\end{align*}
	\noindent уже известно, что $x_1 = -1$, значит из первого уравнения $x_3$ линейно выражается через $x_2$.\\
	из второго $x_4$ линейно выражается через $x_2$ и $x_3$ то есть линейно через $x_2$\\
	в конце концов все $x_k, \ k \in 1 \ldots m $ линейно выражаются через $x_2$\\
	Для окончательного нахождения плана достаточно найти $x_2$ 
	\textbf{ИЛИ} найти $a$ и $b$.\\
	
	
	\noindent Заметим, что узлы удовлетворяют
	\[
	p_0'(x) = 
	\frac{d}{dt} \left[\prod_{k=1}^{m} (t - x_k) \right]_{t=x} = 
	\frac{d}{dt} \left[\prod_{k=1}^{m} (t - (at_{k-1} + b) ) \right]_{t=x} = 
	 0
	\]
	
	\noindent Введём новую переменную $s = \frac{t-b}{a}$
	\[
	\frac{d}{dt} \left[\prod_{k=1}^{m} (t - (at_{k-1} + b) ) \right] =
	\frac{1}{a}\frac{d}{dt} \left[\prod_{k=1}^{m} \left( \frac{t-b}{a} - t_{k-1} \right) \right] 
	\]
	\[
	\frac{d}{dt} \left[\prod_{k=1}^{m} \left( \frac{t-b}{a} - t_{k-1} \right) \right] = 0 
	\quad \Longleftrightarrow \quad
	\frac{d}{ds} \left[\prod_{k=1}^{m} \left( s - t_{k-1} \right) \right] = 0
	\]
	
	\noindent В $x_k = a t_{k-1} + b, k \in 1, \ldots m$ подставляем $k=1$ и получаем
	\[
	x_1 = a t_0 + b
	\]
	\[
	-1 = a (-1) + b
	\]
	\[
	a = b + 1
	\]
	
	\noindent Когда мы найдем значения $s$ из
	\[
	\frac{d}{ds} \left[\prod_{k=1}^{m} \left( s - t_{k-1} \right) \right] = 
	\frac{d}{ds} \left[\prod_{k=0}^{m-1} \left( s - t_{k} \right) \right] = 0
	\]
	
	\noindent Мы воспользуемся тем, что
	\[
	s |_{t=x} = \frac{x-b}{a} = \frac{x-b}{b+1}
	\]
	
	\noindent И отсюда найдём $b$, затем $a$ \\
	после этого можно явно выписать выражения для $x_k \  k \in 1 \ldots m$
	
	\noindent Если обратиться к полиномам Чебышёва 2-го рода то можно переписать уравнение на $s$ в виде:
	\[
	\frac{d}{ds} \left[ (s-t_0) \prod_{k=1}^{m-1} \left( s - t_{k} \right) \right] = 
	\frac{d}{ds} \left[ (s+1) \prod_{k=1}^{m-1} \left( s - t_{k} \right) \right] = 
	C \frac{d}{ds} \left[ (s+1) U_{m-1}(s) \right] = 0
	\]
	
	\noindent Где
	\[
	U_{m-1}(s) = \frac{\sin(m \arccos(s))}{\sin(\arccos(s))}
	\]
	
	\subsection{Ограничение на $x$ от узлов}
	
	\noindent Вспоминаем геометрию узлов:
	\[
	-1 = x_1 < x_2 < \ldots < x_m < 1 < z
	\]
	
	\noindent Отсюда, с учетом растяжения вправо ограничение на $x$:
	\[
	x_m = x_m(x) \in (t_{m-1}, \ 1)
	\]
	
	\newpage
	\subsection*{Если растяжка влево}
	
	\noindent В таком сетапе считаем
	\[
	z < -1 < x_1 < \ldots < x_{m-1} < x_m = 1
	\]
	
	\noindent где $x_k = a t_{k} + b, k \in 1, \ldots m$ \\
	и $z=a t_0 + b$, $a > 0$ (её игнорим)\\
	для неких $a, b \in \mathbb{R}$ (мы их не знаем, но они явно и не нужны, главное что какое-то линейное преобразование) \\
	а $t_k, \  k \in 0, \ldots, m$ есть точки альтернанса полинома Чебышёва 1-го рода 
	\[
	T_{m}(x) = \cos( m \arccos(x))
	\]
	\noindent Они имеют вид 
	\[
	t_k = \cos \left( \frac{ (m-k) \pi}{m} \right), \ k \in 0 \ldots m
	\]
	
	\noindent Уравнения описывающие растяжение (отношение длин интервалов между точками альтернанса сохраняется)
	\begin{align*}
		\frac{x_{m-1} - x_{m-2}}{x_m - x_{m-1}} &= \frac{t_{m-1} - t_{m-2}}{t_m - t_{m-1}}\\[0.1in]
		\frac{x_{m-2} - x_{m-3}}{x_m - x_{m-1}} &= \frac{t_{m-1} - t_{m-2}}{t_m - t_{m-1}}\\[0.1in]
		&\cdots\\[0.1in]
		\frac{x_2 - x_1}{x_m - x_{m-1}} &= \frac{t_2 - t_1}{t_m - t_{m-1}}\\
	\end{align*}
	
	\noindent уже известно, что $x_m = -1$, значит из первого уравнения $x_{m-2}$ линейно выражается через $x_{m-1}$.\\
	из второго $x_{m-3}$ линейно выражается через $x_{m-2}$ и $x_{m-1}$ то есть линейно через $x_{m-1}$\\
	в конце концов все $x_k, \ k \in 1 \ldots m $ линейно выражаются через $x_{m-1}$\\
	Для окончательного нахождения плана достаточно найти $x_{m-1}$ \textbf{ИЛИ} найти $a$ и $b$.\\
	
	
	\noindent Для поиска узлов используем
	\[
	p_0'(x) = 
	\frac{d}{dt} \left[\prod_{k=1}^{m} (t - x_k) \right]_{t=x} = 
	\frac{d}{dt} \left[\prod_{k=1}^{m} (t - (at_{k} + b) ) \right]_{t=x} = 
	0
	\]
	
	\noindent Введём новую переменную $s = \frac{t-b}{a}$
	\[
	\frac{d}{dt} \left[\prod_{k=1}^{m} (t - (at_{k} + b) ) \right] =
	\frac{1}{a}\frac{d}{dt} \left[\prod_{k=1}^{m} \left( \frac{t-b}{a} - t_{k} \right) \right] 
	\]
	\[
	\frac{d}{dt} \left[\prod_{k=1}^{m} \left( \frac{t-b}{a} - t_{k} \right) \right] = 0 
	\quad \Longleftrightarrow \quad
	\frac{d}{ds} \left[\prod_{k=1}^{m} \left( s - t_{k} \right) \right] = 0
	\]
	
	\noindent В $x_k = a t_{k} + b, k \in 1, \ldots m$ подставляем $k=m$ и получаем
	\[
	x_m = a t_m + b
	\]
	\[
	1 = a (1) + b
	\]
	\[
	a = 1 - b
	\]
	
	\noindent Когда мы найдем значения $s$ из
	\[
	\frac{d}{ds} \left[\prod_{k=1}^{m} \left( s - t_{k} \right) \right] = 0
	\]
	
	\noindent Мы воспользуемся тем, что
	\[
	s |_{t=x} = \frac{x-b}{a} = \frac{x-b}{1 - b}
	\]
	
	\noindent И отсюда найдём $b$, затем $a$ \\
	после этого можно явно выписать выражения для $x_k \  k \in 1 \ldots m$
	
	
	\noindent Если обратиться к полиномам Чебышёва 2-го рода то можно переписать уравнение на $s$ в виде:
	\[
	\frac{d}{ds} \left[ (s-t_m) \prod_{k=1}^{m-1} \left( s - t_{k} \right) \right] = 
	\frac{d}{ds} \left[ (s-1) \prod_{k=1}^{m-1} \left( s - t_{k} \right) \right] = 
	C \frac{d}{ds} \left[ (s-1) U_{m-1}(s) \right] = 0
	\]
	
	\noindent Где
	\[
	U_{m-1}(s) = \frac{\sin(m \arccos(s))}{\sin(\arccos(s))}
	\]
	
	
	\subsection{Ограничение на $x$ от узлов}
	
	\noindent Вспоминаем геометрию узлов:
	\[
	z < -1 < x_1 < \ldots < x_{m-1} < x_m = 1
	\]
	
	\noindent Отсюда, с учетом растяжения влево ограничение на $x$:
	\[
	x_1 = x_1(x) \in (-1, \ t_1)
	\]
	
	
	\newpage
	\section{Два узла на границе}
	
	\subsection{Веса}
	
	\noindent Они получаются такими же как и в получебышёвском случае:
	\[
	\omega_k  = 
	\frac
	{
		(-1)^{k+1} L_k'(x)
	}
	{
		\sum_{k=1}^{m} (-1)^{k+1}  L_k'(x)
	}, \ \forall k \in 1 \ldots m 
	\]
	
	\[
	L_k(x) = \underset{k \neq j}{\prod_{j=1}^{m}} \frac{x - x_j}{x_k - x_j}, \ k \in 1 \ldots m
	\]
	
	\subsection{Ограничение на $x$ от весов}
	
	Мы требуем $\forall k \in 1 \ldots m \ \omega_k > 0 $
	
	\noindent То есть нужно решить систему неравенств относительно $x$:
	\[
	\omega_k  = 
	\frac
	{
		(-1)^{k+1} L_k'(x)
	}
	{
		\sum_{k=1}^{m} (-1)^{k+1}  L_k'(x)
	} > 0 , \ \forall k \in 1 \ldots m 
	\]
	
	\subsection{Узлы}
	
	Теперь узлы должны быть расположены следующим образом:
	\[
	-1 = x_1 < x_2 < \ldots < x_{m-1} < x_m = 1
	\]
	
\end{document}